% !TeX root = ../../main.tex
% The three stages of scintillation in an activated inorganic crystal.
% \input by chapters/02-literature-review.tex.
%
% The two bands and the dividers at x = 5.2 and 10.4 are the frame; everything
% else hangs off a node inside one stage, so an edit moves a group rather than
% orphaning a label from what it describes.
%
% THREE INVARIANTS KEEP THE LABELS APART. Hold them and nothing can collide;
% break one and TeX will not say so.
%   1. Band labels are pinned to the east end of each band, so the middle of
%      both bands belongs to the carriers and no dot can land on the text.
%   2. Every prose block has a text width, so its box is fixed rather than set
%      by however long the string happens to be. Stage 1 starts at x = 2.25 and
%      is 2.8 cm wide, which is what holds it between track and divider.
%   3. In stage 3 the label column's LEFT EDGE IS THE ACTIVATOR'S west edge and
%      both carrier arrows terminate on that same edge, so every arrow is
%      always left of every label. Anchor the labels anywhere else and that
%      guarantee is gone.
% \scintfit checks the one thing the invariants cannot: total width.
%
% Needs arrows.meta and decorations.pathmorphing, both loaded in main.tex.

\providecolor{scintcbfill}{HTML}{E6F1FB}
\providecolor{scintcbline}{HTML}{378ADD}
\providecolor{scintvbfill}{HTML}{FBEAF0}
\providecolor{scintvbline}{HTML}{D4537E}
\providecolor{scintactfill}{HTML}{FAEEDA}
\providecolor{scintactline}{HTML}{EF9F27}
\providecolor{scinthot}{HTML}{BA7517}
\providecolor{scintgamma}{HTML}{E24B4A}
\providecolor{scintinert}{HTML}{888780}

% Width guard, called just before \end{tikzpicture}. The figure is drawn at the
% full measure, so a label added at either end overruns the text block in
% silence; this raises a hard error instead.
\providecommand{\scintfit}{%
	\path let \p1=(current bounding box.east), \p2=(current bounding box.west) in
		\pgfextra{%
			\ifdim\dimexpr\x1-\x2\relax>\textwidth
				\PackageError{scintillation-process}{Figure is wider than the text block}%
					{It measures \the\dimexpr\x1-\x2\relax\space against a measure of
					\the\textwidth. Shorten a label or move it inside the diagram.}%
			\fi};}

\begin{tikzpicture}[
	x=1cm, y=1cm,
	ann/.style    = {font=\footnotesize, align=center},
	block/.style  = {ann, align=left, text width=2.8cm},
	% 3.8 cm: the widest single line in the stack measures 3.67 cm, and below
	% that it wraps and the stack grows into the valence band
	stack/.style  = {ann, text width=3.8cm, inner ysep=1pt},
	title/.style  = {font=\small},
	e/.style      = {circle, fill=scintcbline, inner sep=0pt, minimum size=2.4pt},
	h/.style      = {circle, fill=scintvbline, inner sep=0pt, minimum size=2.4pt},
	hot/.style    = {circle, fill=scintactline, inner sep=0pt, minimum size=2.4pt},
	flow/.style   = {-{Stealth[length=4pt]}, line width=0.4pt},
	actbox/.style = {draw=scintactline, fill=scintactfill, line width=0.4pt,
		rounded corners=1pt, minimum width=2cm, minimum height=0.5cm,
		inner sep=0pt, font=\footnotesize, text=scintactline!55!black},
	wavy/.style   = {-{Stealth[length=4pt]}, line width=0.5pt, decorate,
		decoration={snake, amplitude=0.5mm, segment length=2mm,
			post length=1.6mm}},
	]

% ---------- frame: bands and stage dividers --------------------
\def\scintR{15.4}

\draw[dashed, scintinert, line width=0.3pt] (5.2,0.1) -- (5.2,7.8);
\draw[dashed, scintinert, line width=0.3pt] (10.4,0.1) -- (10.4,7.8);

\fill[scintcbfill, draw=scintcbline, line width=0.4pt, rounded corners=1pt]
	(0.6,6.0) rectangle (\scintR,6.4);
\fill[scintvbfill, draw=scintvbline, line width=0.4pt, rounded corners=1pt]
	(0.6,1.6) rectangle (\scintR,2.0);

% invariant 1: pinned east, leaving the middle of both bands to the carriers
\node[font=\footnotesize, scintcbline!60!black, anchor=east]
	at ([xshift=-2mm]\scintR,6.2) {Conduction band};
\node[font=\footnotesize, scintvbline!60!black, anchor=east]
	at ([xshift=-2mm]\scintR,1.8) {Valence band};

\draw[{Stealth[length=4pt]}-, line width=0.3pt] (0.28,2.0) -- (0.28,3.55);
\draw[-{Stealth[length=4pt]}, line width=0.3pt] (0.28,4.45) -- (0.28,6.0);
\node[font=\footnotesize] at (0.28,4.0) {$E_\mathrm{g}$};

% ---------- stage headings -------------------------------------
\node[title]	at (2.9,8.65)  {1.\ Conversion};
\node[ann]		at (2.9,8.15)  {$\sim\!10^{-15}\,$s};
\node[title]	at (7.8,8.65)  {2.\ Transport};
\node[ann]		at (7.8,8.15)  {$10^{-12}$--$10^{-10}\,$s};
\node[title]	at (12.9,8.65) {3.\ Luminescence};
% Ce:LSO decays in ~40 ns and NaI(Tl) in ~230 ns, so the range has to cover
% both activators named below, not just the fast one.
\node[ann]		at (12.9,8.15) {$10^{-8}$--$10^{-7}\,$s decay};

% ================ STAGE 1: conversion ==========================
\draw[wavy, scintgamma] (0.75,0.55) -- (1.88,1.53);
\node[ann, anchor=west] at (0.6,0.2) {\keV{511} photon};
\draw[-{Stealth[length=4pt]}, scinthot, line width=0.4pt] (1.95,1.60) -- (2.05,7.40);

% cascade: {parent}/{child} coordinate pairs -- edit branch geometry here
\foreach \a/\b in {
	{2.05,7.40}/{2.40,7.05},
	{2.40,7.05}/{2.25,6.55},
	{2.40,7.05}/{2.75,6.85},
	{2.75,6.85}/{2.65,6.48},
	{2.75,6.85}/{3.10,7.10},
	{3.10,7.10}/{3.30,6.62},
	{3.10,7.10}/{3.45,7.30},
	{3.45,7.30}/{3.75,6.90},
	{3.75,6.90}/{4.00,6.50}}
	\draw[scintactline, line width=0.4pt] (\a) -- (\b);

\foreach \p in {{2.25,6.55},{2.65,6.48},{3.30,6.62},{4.00,6.50},{3.45,7.30}}
	\node[hot] at (\p) {};
\foreach \x in {2.25,4.00}
	\draw[flow, scintinert] (\x,6.44) -- (\x,6.30);

\foreach \x in {2.0,2.6,3.2} \node[h] at (\x,1.80) {};

% invariant 2: a fixed 2.8 cm measure from x = 2.25 clears the track on the
% left and the divider on the right, whatever the wording becomes
\node[block, anchor=west] (s1cascade) at (2.25,4.95)
	{Impact ionisation and Auger cascade};
\node[block, anchor=north west] (s1thermal)
	at ([yshift=-4mm]s1cascade.south west)
	{Carriers thermalise to the band edges};

% ================ STAGE 2: transport ===========================
\foreach \x in {5.8,7.0,8.2} \node[e] at (\x,6.20) {};
\foreach \x in {5.8,7.0,8.2} \node[h] at (\x,1.80) {};
\draw[flow, scintcbline] (5.8,5.60) -- (9.70,5.60);
\draw[flow, scintvbline] (5.8,2.50) -- (9.70,2.50);

\draw[dashed, scintinert, line width=0.4pt] (5.80,4.50) -- (7.00,4.50);
\node[e] at (6.40,4.62) {};
\draw[flow, scintinert] (6.00,5.35) -- (6.30,4.72);
\node[ann] at (6.40,4.05) {Trapping};

\node[e] at (8.90,4.50) {};
\node[h] at (9.15,4.50) {};
\draw[dashed, scintinert, line width=0.3pt] (9.02,4.50) ellipse (0.45 and 0.30);
\node[ann] at (9.02,4.05) {Exciton};

\node[ann, text width=4.0cm] at (7.8,3.05)
	{Diffusion down the\\[-1pt] concentration gradient ---\\[-1pt]
		no applied field};

% ================ STAGE 3: luminescence ========================
% invariant 3: the label column starts at activator.west, and both arrows
% terminate on that same edge, so no arrow can reach a label
\node[actbox] (activator) at (12.90,4.45) {Activator};

\draw[flow, scintcbline] (11.00,6.00) -- (activator.north west);
\draw[flow, scintvbline] (11.00,2.00) -- (activator.south west);

\node[stack, anchor=south west] (recomb)
	at ([yshift=4mm]activator.north west) {Radiative recombination};
\node[stack, anchor=north west] (dopants)
	at ([yshift=-3mm]activator.south west)
	{Ce$^{3+}$ (LSO, LYSO)\\[-1pt] or Tl$^{+}$ (NaI)};
\node[stack, anchor=north west] (emitted)
	at ([yshift=-2mm]dopants.south west) {Optical photons, $h\nu < E_\mathrm{g}$};

% fanned into the gap the stack leaves between recomb and dopants
\draw[wavy, scinthot] ([yshift=1.2mm]activator.east) -- ++(1.25,0.35);
\draw[wavy, scinthot] (activator.east) -- ++(1.35,0);
\draw[wavy, scinthot] ([yshift=-1.2mm]activator.east) -- ++(1.25,-0.25);

% ---------- inset: track confinement ---------------------------
\fill[scintinert!12, draw=scintinert, line width=0.3pt, rounded corners=1pt]
	(0.6,-1.35) rectangle (2.4,-0.35);
\draw[scintactline, line width=1.4pt, opacity=0.75, line cap=round]
	(1.40,-1.25) -- (1.55,-0.45);
\node[ann, anchor=west, align=left] at (2.7,-0.85)
	{Free carriers exist only inside a $\sim$\um{1} ionisation track,\\[-1pt]
		for picoseconds --- the surrounding bulk stays insulating};

% ---------- legend ---------------------------------------------
\node[e] at (0.75,-1.95) {};
\node[ann, anchor=west] at (0.95,-1.95) {Electron};
\node[h] at (3.00,-1.95) {};
\node[ann, anchor=west] at (3.20,-1.95) {Hole};
\node[hot] at (4.85,-1.95) {};
\node[ann, anchor=west] at (5.05,-1.95) {Hot carrier};

\scintfit
\end{tikzpicture}
